\chapter{Introduction}
\label{chap:introduction}

\section{Context and Motivation}
\label{sec:intro_motivation}

\par Intrusion Detection Systems (IDS) are a core defensive layer in modern networks, monitoring traffic to detect malicious behaviour such as denial-of-service attacks, reconnaissance scans, brute-force attempts, and data exfiltration. The classical approach is signature-based detection, in which traffic is matched against a database of hand-written rules describing known attacks; the widely deployed Snort system is the canonical example \cite{Roesch1999}. Signature matching is fast and precise on threats it already knows, but it has two structural weaknesses: it cannot recognise an attack for which no rule has been written, and the rule set must be maintained by hand as new threats appear.

\par These weaknesses have become more pressing with the growth of the Internet of Things (IoT). Modern deployments place large numbers of constrained, infrequently patched devices on the same networks as critical services, and the volume and diversity of the resulting traffic strain a detection model that depends on enumerating every attack in advance. The same devices are attractive targets: compromised IoT endpoints are routinely conscripted into large botnets and used to launch volumetric denial-of-service campaigns, which is reflected in benchmark data where denial-of-service and distributed denial-of-service flows dominate the attack distribution \cite{CICIoT2023}.

\par Machine learning, and deep learning in particular, offers an alternative that learns the statistical structure of traffic directly from data rather than from hand-written rules \cite{Ferrag2020, Lansky2021}. A model trained on labelled flows can, in principle, generalise to variants of known attacks and flag anomalous behaviour without an exact signature. This premise has produced a large literature reporting strong benchmark accuracy for neural intrusion detection, including work that applies deep models to real-time web intrusion detection rather than offline analysis alone \cite{AIIDS2020, Wu2022}.

\section{The Research-to-Deployment Gap}
\label{sec:intro_gap}

\par Strong benchmark numbers do not by themselves yield a deployable detector. This thesis addresses the gap between the two \cite{Apruzzese2023}. Three issues recur. First, many reported results are obtained on random train/test splits, which let near-duplicate flows from the same capture session appear in both partitions and inflate the measured accuracy relative to what the model would achieve on genuinely unseen traffic. Second, network traffic drifts: attack tooling and benign usage both change over time, so a model evaluated without respect to temporal order is not tested under the conditions it will actually face \cite{Andresini2021}. Third, a deployed detector must do more than emit a label: it must report a confidence that can be trusted, and modern neural networks are known to be systematically overconfident unless their outputs are explicitly calibrated \cite{Guo2017}. To these we add the practical constraint that real-time operation imposes a latency budget an offline benchmark never measures.

\par This thesis builds these concerns into the design from the outset. Evaluation uses a temporally ordered split so that the test set is strictly later than the training data; the deep model is measured against strong tabular baselines under identical conditions so that any claimed advantage is real; the classifier's confidence is calibrated before it is shown to a user; and inference latency is benchmarked as an explicit acceptance criterion.

\section{Approach and Scope}
\label{sec:intro_approach}

\par The work is built on the CICIoT2023 dataset, a large-scale IoT intrusion-detection benchmark of over 46 million labelled flow records spanning 33 attack types, captured in a 105-device smart-home testbed \cite{CICIoT2023}. The primary classifier is a Multi-Layer Perceptron (MLP) trained at two granularities: a binary attack-versus-benign task and an eight-class task over coarse attack families. Because the public CICIoT2023 features are tabular flow statistics rather than raw packets, the deep model is not assumed to be the best tool by default; recent evidence shows that well-tuned simple neural networks are competitive with gradient-boosted trees on tabular data \cite{Kadra2021}, and the thesis tests this empirically by comparing the MLP against Random Forest and XGBoost baselines on the same splits. Beyond offline evaluation, the trained model is deployed behind a live demonstration in which scripted attacks are launched against a controlled target and the resulting flows are classified in near real time.

\par The scope is bounded in two respects. The system is trained and evaluated only on the 39-feature public release of CICIoT2023, which is derived from packet headers and flow metadata and carries no payload features; detection of application-layer attacks, whose signal lies in the payload, is therefore expected to be limited, and the results chapter reports this rather than working around it. Generalisation beyond the smart-home setting of the dataset is not claimed.

\section{Objectives and Contributions}
\label{sec:intro_contributions}

\par The thesis makes six contributions. The first is a reproducible end-to-end pipeline from raw per-folder CICIoT2023 captures to a live per-flow classifier, with deterministic seeding and a fixed artefact contract so that results can be regenerated and the trained model served without re-running training. The second is a temporal evaluation of an MLP against Random Forest and XGBoost baselines at binary and eight-class granularities, reporting accuracy, macro-F1, and weighted-F1 on the train, validation, and test partitions. The third is a feature analysis that identifies the most discriminative flow features and prunes the 39 public features to a 25-feature set, dropping near-zero-variance flags and redundant continuous columns found through variance and correlation analysis. The fourth is confidence calibration of the deep model by temperature scaling \cite{Guo2017}, which lowers the expected calibration error so that the confidence shown to a user reflects empirical accuracy instead of sitting uniformly near certainty. The fifth is a deployed demonstration system, NEURAL IDS, comprising a containerised inference backend and a web dashboard, which validates the core classification functionality against both benign and synthetic-attack traffic. The sixth is an account of where the approach falls short, particularly the eight-class attack-family task, on which no model reaches the 0.85 weighted-F1 bar met on the binary task; this shortfall stems from the fact that the public features describe transport-layer behaviour only.

\section{Thesis Structure}
\label{sec:intro_structure}

\par The remainder of this thesis is organised as follows. Chapter~\ref{chap:theoretical_background} presents the theoretical background, covering intrusion detection systems, deep learning for traffic classification, and the CICIoT2023 dataset. Chapter~\ref{chap:requirements} states the application requirements and specification, defining the problem, the actors involved, and the functional and non-functional constraints the system must satisfy. Chapter~\ref{chap:design_implementation} describes the design, implementation, and testing of the application, including the data pipeline, the model architecture and training procedure, the live-demo environment, and the verification of the core functionality. Chapter~\ref{chap:results} reports the experimental results and discusses the findings. Chapter~\ref{chap:conclusions} concludes the thesis and outlines directions for future work.

\section{Declaration on the Use of Generative AI Tools}
\label{sec:intro_ai_declaration}

\par Generative AI tools were used during the preparation of this thesis. AI-based coding assistants supported the development of the training pipeline and the demonstration application by generating boilerplate code, suggesting refactorings, and helping diagnose errors. AI language models were also used to draft and revise portions of the written text for clarity and concision. All such output was reviewed, tested, and edited by the author. Every experimental result reported here was produced by running the pipeline and was verified independently of any AI tool, and the author takes full responsibility for the content of the thesis, the correctness of the results, and the final wording.
